% Corte mecánico íntegro: parte_ii.tex:323--419.
% Residencia material del ensamblaje sucesor de 102 capítulos.
% Residencia causal: capítulo 46.

\chapter{Spectral Pell screw and logarithmic suspension}

\section{Screw law}

The corpus contains the exact identity
\begin{equation}
 V^{-1}\mathscr S V=(2+\sqrt3)\ii\,\mathscr S.
 \label{espiral:eq:tornillo-pell}
\end{equation}
Let \(\varepsilon=2+\sqrt3\). The orbit transversal
\[
 z_n=z_0(\varepsilon\ii)^n
\]
satisfies
\[
 r_n=r_0\varepsilon^n,
 \qquad
 \theta_n=\theta_0+\frac{\pi n}{2}.
\]
Eliminando \(n\),
\begin{equation}
 r(\theta)=r_0
 \exp\!\left[
 \frac{2\log(2+\sqrt3)}{\pi}
 (\theta-\theta_0)
 \right].
 \label{espiral:eq:espiral-log-pell}
\end{equation}
The logarithmic spiral is an output of the operator: its equation was not chosen
to fit \(\varepsilon\) afterward.

\begin{proof}[Derivation of \cref{espiral:eq:espiral-log-pell}]
From \(\theta_n-\theta_0=n\pi/2\) it follows that
\(n=2(\theta_n-\theta_0)/\pi\). Substituting into
\(r_n=r_0\varepsilon^n\),
\[
 r_n=r_0\exp\left(
 \frac{2\log\varepsilon}{\pi}(\theta_n-\theta_0)
 \right).
\]
Continuous interpolation is the indicated equation.
\end{proof}

\section{Conjugate contraction involution}

The bilateral extension \(n\in\Z\) produces
\[
 n\to+\infty:\ r_n\to\infty,
 \qquad
 n\to-\infty:\ r_n\to0.
\]
The inversion \(n\mapsto-n\) interchanges expansion and contraction. Since
\(\varepsilon^{-1}=2-\sqrt3\), both branches belong to the same Pell unit.
They are not two independent spirals, but two orientations of the same
spectral screw.

\begin{figure}[htbp]
\centering
\begin{tikzpicture}[scale=.82]
 \draw[->,HMTGray] (-5.1,0)--(5.1,0) node[right] {$x$};
 \draw[->,HMTGray] (0,-3.1)--(0,3.1) node[above] {$y$};
 \draw[HMTBlue,very thick,domain=-7:5.2,samples=240,smooth,variable=\t]
   plot ({.27*exp(.22*\t)*cos(deg(\t))},{.27*exp(.22*\t)*sin(deg(\t))});
 \draw[HMTRed,thick,domain=-7:5.2,samples=240,smooth,variable=\t]
   plot ({.27*exp(.22*\t)*cos(deg(-\t))},{.27*exp(.22*\t)*sin(deg(-\t))});
 \foreach \k in {-4,-3,...,4}{
   \pgfmathsetmacro{\ang}{1.57079632679*\k}
   \fill[HMTNavy]
     ({.27*exp(.22*\ang)*cos(deg(\ang))},
      {.27*exp(.22*\ang)*sin(deg(\ang))}) circle (.8pt);
 }
 \node[HMTBlue,font=\small\sffamily] at (3.5,2.4) {expansion};
 \node[HMTRed,font=\small\sffamily] at (3.5,-2.4) {conjugate orientation};
\end{tikzpicture}
\caption{Transverse projection of the Pell screw. The marks represent quarter turns.}
\label{espiral:fig:tornillo-pell}
\end{figure}

\section{Scale covariance and energy conservation}

The dilation \(\varepsilon\) must not be interpreted as gratuitous growth
of the energetic amplitude of a wave in vacuum. Its coherent realization
acts on the spatial or spectral scale:
\[
 (r,\varphi,\lambda)
 \longmapsto
 (\varepsilon r,\varphi+\pi/2,\varepsilon\lambda).
\]
This distinction will be essential in the electromagnetic realization of
\cref{espiral:chap:realizacion-em}.

\begin{resultbox}[title={Exact compatibility between the Pell spectral screw and the phase--memory suspension}]
The minimal phase--memory lift is a circular helix. The spectral screw
combines a quarter-turn and a Pell dilation; its projection is a logarithmic
spiral. Both arise from the HMT ledger, but they are neither the same
periodicity nor obtained from one another by forgetting types.
\end{resultbox}
